\chapterimage{images/19-crkva-i-carstvo.jpg}

\chapter{Crkva i carstvo}

\editorialnote{Novi političko-religijski poredak i njegov uticaj na svakodnevni život i slobodu.}

Paks nije nastao u jednom trenutku.

Rastao je tokom dva i po veka.

Svet po svet.

Generacija po generacija.

Najpre kao vera koja je nudila ono što niko drugi nije mogao.

Povratak iz smrti.

Zatim kao politički savez.

Na kraju kao carstvo koje nije volelo da sebe naziva carstvom.

\bigskip

Njegov centar bio je Pacem.

Vatikan.

Papa.

Rimska kurija.

Svete kongregacije.

Ispod njih protezala se mreža biskupa, kardinala, sveštenika, lokalnih vlada, vojske, trgovine i administracije preko stotina svetova.

\bigskip

Formalno, mnoge planete još su imale svoje vlade.

Gradonačelnike.

Savete.

Kraljeve.

Parlamente.

Ali postojalo je pitanje iznad svih tih institucija.

Da li pripadaju Paksu?

\bigskip

Ako pripadaju, onda Crkva nije bila samo još jedna religija.

Bila je izvor autoriteta.

\bigskip

Papsko naređenje moglo je da nadjača admirale, generale i civilne upravnike.

De Soja je to video neposredno.

Bio je samo kapetan jednog broda.

Ali kada mu je Papa dao posebna ovlašćenja za potragu za Aneom, admirali sa flotama morali su da ga slušaju.

Civilni lideri morali su da mu otvore svoje gradove.

Lokalni biskupi morali su da mu daju ljude i sredstva.

\bigskip

De Soja nikada nije sebe smatrao političarem.

Bio je jezuita.

Sveštenik.

Vojnik.

Verovao je da služi Bogu i štiti ljude.

Upravo zato njegova perspektiva nije bila jednostavna.

Paks za njega nije bio samo tiranija.

Video je šta je doneo mnogim svetovima.

Red posle vekova haosa.

Trgovinu.

Medicinu.

Odbranu.

Veze između udaljenih planeta.

I vaskrsenje.

\bigskip

De Soja je sam umirao.

Više puta.

\bigskip

Kurirski brod \emph{Rafael} omogućavao je Paksu nešto što Hegemonija nikada nije mogla da uradi.

Posada bi legla u komore.

Brod bi ubrzao silom koju živo ljudsko telo ne može da podnese.

Tela bi bila uništena.

Brod bi stigao do drugog sistema mnogo brže nego obična letelica.

Tamo bi kruciforma i medicinski sistemi ponovo izgradili putnike.

\bigskip

Za De Soju, buđenje posle smrti nije bilo samo tehnologija.

Bilo je religijsko iskustvo.

Svaka ćelija bolela ga je.

Sećanje na smrt bilo je stvarno.

Ali jednako stvaran bio je osećaj da mu je život ponovo darovan.

\bigskip

Zato je razumeo zašto ljudi veruju.

Ako čovek jednom umre i zatim otvori oči, teško je posmatrati kruciformu kao običan medicinski uređaj.

\bigskip

Papa Julije bio je najveći dokaz.

Sveti Otac vladao je gotovo dva i po veka.

Umirao je.

Vraćao se.

Ponovo bio krunisan.

I ponovo.

\bigskip

Njegovo prvo ime bilo je Lenar Hojt.

\bigskip

Hodočasnik sa Hiperiona.

Sveštenik koji je nekada u sebi nosio dve kruciforme.

Čovek čija je smrt vratila Pola Direa.

\bigskip

Ali istorija Crkve posle Pada bila je mnogo komplikovanija od zvanične verzije.

Dire je jedno vreme postao Papa.

Tejar I.

Njegova Crkva trebalo je da bude otvorena prema zamisli da čovečanstvo još nije završena vrsta.

Da se svest menja.

Razvija.

Možda čak približava nečemu što ljudi nazivaju božanskim.

\bigskip

Posle Direove smrti ponovo se vratio Hojt.

Kao Julije VI.

I promenio smer Crkve.

\bigskip

Tejarova učenja proglašena su opasnim.

Kasnije jeretičkim.

Redovi koji su ih zadržali raspušteni su.

Sveštenici koji nisu prihvatili novu doktrinu uklanjani su iz centara moći.

Neki su poslati na zabačene svetove.

\bigskip

Otac Glaukus bio je jedan od njih.

Njegova kazna bila je Sol Drakoni Septem.

Ledeni svet na kraju ničega.

Crkva ga nije ubila.

Samo ga je poslala tamo gde njegov glas više nije mogao da smeta.

\bigskip

Takva je često bila priroda Paksa.

Nije svaki protivnik završavao pred streljačkim vodom.

Nekada je bilo dovoljno oduzeti mu publiku.

Položaj.

Putovanje.

Pristup obrazovanju.

Mogućnost da govori dovoljno glasno da ga drugi čuju.

\bigskip

Kruciforma je sve to činila mnogo lakšim.

\bigskip

Preobraćenje nije bilo samo pitanje spasenja.

Otvaralo je vrata.

\bigskip

Verujući su lakše dobijali položaje.

Mogli su da putuju brodovima Paksa.

Da koriste sistem vaskrsavanja.

Da postanu deo administracije koja povezuje svetove.

Na mnogim planetama najbogatiji i najuticajniji slojevi prvi su prihvatali novu veru.

Posle njih dolazili su ostali.

\bigskip

Oni koji su odbijali kruciformu nisu nužno bili zatvarani.

Ali živeli su u svetu koji je sve manje bio projektovan za njih.

\bigskip

Rol je to znao od detinjstva.

Bez kruciforme nikada ne bi mogao ozbiljno da napreduje u vojsci Paksa.

Međuzvezdano putovanje bilo je mnogo teže.

Mnogi poslovi bili su mu zatvoreni.

Njegova odluka da ostane nekršten nije izgledala kao pobuna.

Ali u društvu u kojem gotovo svako prihvata isti znak na grudima, odbijanje je samo po sebi postajalo političko.

\bigskip

Na nekim svetovima prelazak u Paks bio je miran.

Na drugima nije.

\bigskip

Barnardov Svet bio je dobar primer.

Nekada je bio poznat po univerzitetima, malim gradovima i poljoprivredi.

Posle Pada mogao je da preživi sam.

Zato nije imao mnogo razloga da prihvati vlast sa Pacema.

\bigskip

Paks je stigao nekoliko decenija kasnije.

Otpor je trajao generacijama.

Na kraju je završio građanskim ratom između katolika i grupa koje su sebe nazivale Slobodnim Vernicima.

Tek više od dva veka posle Pada Barnardov Svet formalno je uključen u Paks.

\bigskip

Kada je De Soja stigao tamo, rat je već bio istorija.

Polja su bila mirna.

Gradovi uređeni.

Crkve pune.

\bigskip

Ali mnogi stari koledži bili su prazni.

Drugi su pretvoreni u semeništa.

U šumama su još postojale male grupe otpora.

\bigskip

Za nekoga ko bi došao samo na nekoliko dana, svet je izgledao spokojno.

Za nekoga ko poznaje njegovu istoriju, taj mir imao je cenu.

\bigskip

Na drugim planetama cena je bila još vidljivija.

Lokalne kulture preuređivane su tako da odgovaraju novom poretku.

Religije koje nisu prihvatale autoritet Vatikana gurane su na marginu.

Planete koje su odbijale pristupanje nalazile su se izvan zaštite i trgovine Paksa.

\bigskip

Postojala je i vojska.

\bigskip

Flota Paksa.

Kopnene snage.

Hristovi Legionari.

Vatikanska Švajcarska garda.

\bigskip

Zvanično, postojali su da brane ljudske svetove od Proteranih.

I opasnost Proteranih bila je stvarna.

Granica nije bila mirna.

Borbe su trajale vekovima.

Ali strah od spoljnog neprijatelja davao je Crkvi još jedan razlog za koncentraciju vlasti.

\bigskip

Veliki Zid Paksa čuvao je granicu.

Brodovi su patrolirali hiljadama parseka.

Generacije vojnika učene su da Proterani više nisu pravi ljudi.

Da su promenili ljudsko telo.

Genetiku.

Oblik.

Da su napustili ono što Crkva smatra dostojanstvom čoveka.

\bigskip

De Soja je verovao većini toga.

Ali služba na granici ostavila mu je sumnje.

Video je Proterane.

Video je njihove svetove.

Video je i šta Paks radi u ime odbrane.

\bigskip

Ponekad se pitao da li Proterani možda razumeju nešto o budućnosti čoveka što Paks odbija da razume.

Takva misao za jednog sveštenika bila je opasna.

De Soja bi je odbacio.

Ali vraćala se.

\bigskip

Potera za Aneom učinila je sumnje težim za potiskivanje.

\bigskip

Na početku je dobio jednostavno naređenje.

Pronaći devojčicu.

Dovesti je živu na Pacem.

\bigskip

De Soja ga je prihvatio.

Nije znao zašto je toliko važna.

Pretpostavljao je da Papa zna.

To mu je tada bilo dovoljno.

\bigskip

Ali što je potera duže trajala, naređenja su postajala čudnija.

Čitave flote pomerane su zbog jednog deteta.

Hiljade vojnika pretraživale su kontinente.

Ljudi su ispitivani Istinozborom.

Lokalni biskupi uklanjani su ako bi smetali potrazi.

Samo jedna posebna papska naredba bila je dovoljna da sve druge obaveze postanu nevažne.

\bigskip

De Soja je na Mare Infinitusu mobilisao hiljade ljudi da pronađe trag jednog nepoznatog čoveka.

Kada se lokalni biskup pobunio zbog troškova i zloupotrebe vlasti, De Soja ga je uhapsio.

Tek kasnije je počeo ozbiljno da razmišlja o tome šta je zapravo učinio.

\bigskip

Njegova poslušnost bila je iskrena.

Upravo zato je bila opasna.

\bigskip

A onda je saznao da se misija promenila.

\bigskip

Na Pacemu ga je primio kardinal Simon Avgustino Lurdusami.

Jedan od najmoćnijih ljudi u Crkvi.

De Soji je saopšteno da više nije dovoljno da Aneu zarobi.

\bigskip

Crkva je sada smatrala devojčicu bolešću.

Virusom u Telu Hristovom.

\bigskip

De Soja je nekoliko trenutaka pokušavao da razume šta mu govore.

``Ja treba da pronađem dete.''

``Da.''

``A neko drugi treba da ga ubije.''

``Da.''

\bigskip

To nije bilo naređenje koje je dobio na početku.

Ali njemu nije dato pravo da raspravlja.

\bigskip

Nova ratnica već je bila na njegovom brodu.

Radamant Nemes.

\bigskip

Na prvi pogled nije izgledala naročito neobično.

Sitna žena.

Kratka tamna kosa.

Bleda koža.

Uniforma nove Krstaške Legije.

Kruciforma na grudima.

\bigskip

Lurdusami ju je predstavio kao novu vrstu ratnika.

Ne kao biološki izmenjeno biće.

To bi bilo protiv crkvene doktrine.

Ne kao mašinu.

To bi bilo još gore.

\bigskip

Navodno je bila čovek.

Samo bolje obučena.

Poboljšana.

\bigskip

De Soja joj nije verovao.

\bigskip

Nemes je dobila sopstveni papski disključ.

To je bilo još neobičnije.

Formalno je putovala pod De Sojinom komandom.

Ali za određene zadatke imala je ovlašćenje koje je dolazilo neposredno sa vrha Crkve.

\bigskip

Drugim rečima, De Soja je bio odgovoran za misiju čiji najvažniji deo više nije kontrolisao.

\bigskip

Istovremeno, Crkva je spremala još veći korak.

Krstaški rat.

\bigskip

Lurdusami je De Soji objasnio zvaničnu sliku.

Proterani će jednog dana probiti Veliki Zid.

Uništiće Pacem.

Vatikan.

Paks.

Sve što je čovečanstvo ponovo izgradilo.

\bigskip

Papa je zato odlučio da više ne čeka.

Paks će krenuti u ofanzivu.

\bigskip

Svaki važniji svet dobio je naređenje da uloži ogromne resurse u nove ratne brodove.

Čitava civilizacija počela je da se priprema za sveti rat.

\bigskip

Za vernike to je predstavljeno jednostavno.

Odbrana čovečanstva.

Odbrana vere.

Odbrana prava na večni život.

\bigskip

Ali De Soja je već počinjao da vidi pukotine.

\bigskip

Zašto je Anea toliko važna?

Zašto Crkva više ne želi da je ispita, nego da je ubije?

Zašto se pojavila Nemes?

Zašto mrtvi dalekobacači reke Tetis ponovo rade samo pred devojčicom?

I zašto neko u vrhu Paksa zna mnogo više nego što govori?

\bigskip

Anea je Rolu dala odgovor koji je bio još gori od pitanja.

TehnoCentar nije samo preživeo Pad.

Pomogao je da se izgradi novi poredak.

\bigskip

Kruciforme nisu bile Božji dar.

Njihovo poreklo vodilo je do buduće mašinske Ultimativne Inteligencije.

Poslate su unazad kroz vreme.

Prvi neuspešni oblik testiran je na Bikurama na Hiperionu.

Rezultat su bili ljudi koji su vaskrsavali, ali sa svakim povratkom gubili deo inteligencije i ljudskosti.

\bigskip

Srž je kasnije uklonila te nedostatke.

Usavršila proces.

I dala tehnologiju Crkvi.

\bigskip

Ne Polu Direu.

Lenaru Hojtu.

\bigskip

Papi Juliju.

\bigskip

Rol je neko vreme samo gledao Aneu.

``Hoćeš da kažeš da je Srž Crkvi dala vaskrsenje?''

``Da.''

``Posle svega što se dogodilo?''

``Da.''

``I Crkva je prihvatila?''

Anea je klimnula.

\bigskip

A. Betik je rekao ono što je Rol već mislio.

Politička moć Paksa izrasla je iz kruciforme.

\bigskip

Anea je to nazvala pogodbom.

Crkva je dobila milijarde vernika.

Svetove.

Vojsku.

Carstvo.

I praktičnu pobedu nad smrću.

\bigskip

TehnoCentar je zauzvrat dobio nešto što još nisu potpuno razumeli.

\bigskip

Ali jedno je već bilo očigledno.

Srž nije morala da ponovo izgradi staru Mrežu.

Ljudi su joj napravili novu.

\bigskip

Ne od dalekobacača.

Od kruciformi.

\bigskip

Milijarde ljudi nosile su ih u telima.

Vernici.

Vojnici.

Sveštenici.

Političari.

\bigskip

Paks je smatrao da je kruciforma znak pripadnosti Hristu.

Anea je tvrdila da je istovremeno i deo tehnologije TehnoCentra.

\bigskip

Rol je instinktivno dodirnuo sopstvene grudi.

Tamo nije bilo ničega.

Prvi put otkako je odbio krštenje shvatio je da njegova odluka možda nije bila samo tvrdoglavost.

\bigskip

Anea nije tvrdila da su svi ljudi Crkve zli.

Naprotiv.

Dire je bio dobar čovek.

Glaukus takođe.

De Soja, koliko je mogla da nasluti, možda nije njihov pravi neprijatelj.

Milijarde običnih vernika nisu znale ništa o TehnoCentru.

\bigskip

To je upravo činilo sistem snažnim.

Nije počivao samo na laži.

Počivao je na istinskoj veri.

Istinskoj nadi.

Stvarnom vaskrsenju.

I ljudima koji su verovali da čine dobro.

\bigskip

Srž nije morala da kontroliše svakog čoveka.

Dovoljno je bilo da usmeri instituciju.

\bigskip

De Soja je bio savršen primer.

Pošten.

Saosećajan.

Disciplinovan.

Spreman da rizikuje život za svoje ljude.

\bigskip

I spreman da izvrši gotovo svako naređenje ako veruje da dolazi od Crkve.

\bigskip

Ali granica njegove poslušnosti počela je da se približava.

\bigskip

Nemes ju je učinila vidljivom.

\bigskip

Tokom putovanja do Sol Drakoni Septema De Soja i njegova dva čoveka ležali su mrtvi u komorama.

Brod je usporavao stotinama gravitacija.

Kabina je bila gotovo bez vazduha.

Temperatura ispod nule.

\bigskip

Nemes nije bila mrtva.

\bigskip

Izašla je iz komore.

Hodala brodom pod ubrzanjem koje nijedan čovek ne bi mogao da podnese.

Disala joj nisu trebala.

Hladnoća joj nije smetala.

\bigskip

Dok su ostali čekali vaskrsenje, ona je radila.

Planirala.

Silazila na planetu.

Ubijala.

Vraćala se pre nego što bi iko mogao da vidi da je uopšte napustila brod.

\bigskip

De Soja to još nije znao.

Ali počeo je da primećuje tragove.

Potisnike koji su radili kada nije trebalo.

Delove brodskog dnevnika koji nisu imali smisla.

Praznine u vremenu.

Neobična ponašanja nove ratnice.

\bigskip

Njegova vera u Boga nije slabila.

Počela je da slabi vera u ljude koji tvrde da govore u Božje ime.

\bigskip

To je bila razlika koju Paks nije očekivao od svojih vojnika.

\bigskip

Za Rola je situacija bila jednostavnija.

Nikada nije obećao poslušnost Crkvi.

Anea nije tražila da pređe u novu veru.

Nije mu nudila vaskrsenje.

Nije tvrdila da je nepogrešiva.

\bigskip

Čak je često govorila da ne zna.

\bigskip

Za Rola je to postajalo neobično važna razlika.

Ljudi koji su imali najveću vlast u Paksu gotovo nikada nisu govorili da ne znaju.

\bigskip

Anea je govorila stalno.

\bigskip

Nije znala ko otvara portale.

Nije znala šta Šrajk želi.

Nije znala svaki korak sopstvene budućnosti.

Nije znala zašto baš ona mora da uradi ono što je pred njom.

\bigskip

Ali jednu stvar je znala.

Čovek mora imati pravo da bira.

\bigskip

Pravo da odbije kruciformu.

Pravo da veruje.

Pravo da ne veruje.

Pravo da promeni sopstveno telo.

Pravo da ostane isti.

Pravo da ode.

Pravo da pogreši.

Čak i pravo da umre.

\bigskip

Za Paks je takva ideja predstavljala opasnost.

Jer njegova najveća moć nije počivala na vojnom brodu.

Niti na Švajcarskoj gardi.

Niti na papinom disključu.

\bigskip

Počivala je na jednostavnom pitanju koje je postavljao svakom čoveku.

Zašto bi izabrao smrt kada ti mi možemo dati život?

\bigskip

Većina ljudi nije imala dobar odgovor.

\bigskip

Rol jeste.

``Zato što je moj.''

\bigskip

Anea ga je pogledala.

``Šta?''

``Moj život.''

Dotakao je grudi.

``I moja smrt.''

\bigskip

Devojčica se blago nasmešila.

Nije rekla da je u pravu.

Nije morala.

\bigskip

U isto vreme, daleko od njih, Otac Kapetan de Soja spremao se za još jedno putovanje.

Još jednu smrt.

Još jedno vaskrsenje.

Još jedno papsko naređenje.

\bigskip

Pored njega je putovala Radamant Nemes.

Zvanično ratnica Crkve.

U stvarnosti nešto što je De Soja tek trebalo da razume.

\bigskip

A na Pacemu se pripremao Krstaški rat.

Svetovi su dobijali naređenja da grade brodove.

Vojske su mobilisane.

Propovednici su govorili o odbrani vere.

Milijarde vernika spremale su se da poslušaju.

\bigskip

Paks je izgledao snažnije nego ikada.

Ali ispod njegove vere, vojske i vaskrsenja ležala je ista opasnost koja je nekada uništila Hegemoniju.

Zavisnost od tehnologije čije pravo poreklo ljudi nisu razumeli.

\bigskip

Stara Hegemonija predala je svoju slobodu TehnoCentru zato što je želela trenutna vrata između svetova.

Novi Paks učinio je nešto slično iz mnogo dubljeg razloga.

Nije želeo da se odrekne života.

\bigskip

Anea je predstavljala pretnju jer je počinjala da pokazuje da možda postoji drugačiji put.

Zato Crkva više nije želela samo da je pronađe.

Neko unutar njenog carstva odlučio je da Ona Koja Podučava ne sme odrasti.